\section{Peace promised to people still in danger}

The Last Sunday after Pentecost begins with God's promise of peace and
includes Christ's warning to flee. Paul thanks the Father for a deliverance
already received, while the chants ask to be heard from the depths. After
Communion, the Church still asks for healing. These are the prayers of
people who have received grace and continue to need it.

The Gospel's details resist an easy account of that condition. A person on a
roof must leave possessions behind. A nursing mother faces a danger that
confidence alone cannot remove. False claimants offer signs, yet Jesus
forbids belief in them. His own coming will be public, and his angels will
gather the elect. The promise concerns their belonging to him through
these dangers; the warning requires them to act.

Two interpretations follow. The first asks how an accomplished rescue can
coexist with continuing affliction. Augustine explains why the Gradual
praises God for a rescue still to come. Chrysostom distinguishes the
danger that requires flight from the coming that all will see.
The second asks how
Christ's assurance about prayer can coexist with the prayers for changed
wills, converted hearts and healed faults. Aquinas's account of being
made worthy explains why the Epistle thanks God for the capacity to
act. Theophylact explains how the petitioner entrusts himself to God,
and Bede shows what that trust asks of someone who seeks pardon.

\begin{maptable}
Introit & Jer 29:11, 12, 14; Ps 84(85):2 & The Lord promises peace and return; the psalm verse recalls restoration.\\
Collect & \textit{Excita, quaesumus} & The faithful ask God to arouse their wills for fruitful work and further remedies.\\
Epistle & Col 1:9--14 & Prayer for knowledge, good works and endurance accompanies thanks for rescue into the Son's kingdom.\\
Gradual & Ps 43(44):8--9 & The people praise God for deliverance and promise enduring thanksgiving.\\
Alleluia & Ps 129(130):1--2 & A cry from the depths is sung with Alleluias.\\
Gospel & Mt 24:15--35 & Jesus warns of flight and deception, promises his public coming and gathering, and affirms his abiding words.\\
Offertory & Ps 129(130):1--2 & The cry returns, with its opening repeated, at the offering.\\
Secret & \textit{Propitius esto} & Acceptance of offerings and prayers is sought with the conversion of the offerers' desires.\\
Preface & Most Holy Trinity & Thanksgiving confesses three distinct persons, one divine essence and equal majesty, with the angels' praise.\\
Communion & Mk 11:24 & Christ assures those who pray of God's answer.\\
Postcommunion & \textit{Concede nobis} & Those who have received the mysteries ask that their remaining faults be cured.\\
\end{maptable}

The biblical settings and the distinct questions of composition and event
date appear in the appendix, Scriptural Date and Location
(\ref{sec:date-location}). Psalm numbers here give the Missal's
Vulgate number first, with the common modern number in parentheses.
